\section{选择题}
\begin{enumerate}
	
	\item 设4阶行列式 $D = \begin{vmatrix} 3 & 5 & 7 & -9 \\ 2 & 2 & 2 & 2 \\ a & b & c & d \\ x & y & 8 & 3 \end{vmatrix}$, 则 $A_{11} + A_{12} + A_{13} + A_{14} = ( \ \ \ )$
	\vspace{0.5em}
	\begin{choices}
		\item $0$
		\item $2$
		\item $-2$
		\item 无法确定
	\end{choices}
	
	\item 行列式 $\begin{vmatrix} 2 & 0 & 2 \\ 0 & 1 & -1 \\ 3 & 1 & 0 \end{vmatrix} = ( \ \ \ )$
	\vspace{0.5em}
	\begin{choices}
		\item $6$
		\item $-6$
		\item $-4$
		\item $-8$
	\end{choices}
	
	\item 设 $A = \begin{pmatrix} a_{11} & a_{12} & a_{13} \\ a_{21} & a_{22} & a_{23} \\ a_{31} & a_{32} & a_{33} \end{pmatrix}$, $B = \begin{pmatrix} a_{21}+ka_{23} & a_{22} & a_{23} \\ a_{11}+ka_{13} & a_{12} & a_{13} \\ a_{31}+ka_{33} & a_{32} & a_{33} \end{pmatrix}$, $P_1 = \begin{pmatrix} 0 & 1 & 0 \\ 1 & 0 & 0 \\ 0 & 0 & 1 \end{pmatrix}$, $P_2 = \begin{pmatrix} 0 & 1 & 0 \\ 0 & 0 & 1 \\ 1 & 0 & 0 \end{pmatrix}$, $P_3 = \begin{pmatrix} 1 & 0 & k \\ 0 & 1 & 0 \\ 0 & 0 & 1 \end{pmatrix}$, $P_4 = \begin{pmatrix} 1 & 0 & 0 \\ 0 & 1 & 0 \\ k & 0 & 1 \end{pmatrix}$.
	
	下列能够将 $A$ 变换为 $B$的是 ( \ \ \ )
	\vspace{0.5em}
	\begin{choices}
		\item $P_1 A P_3 = B$
		\item $P_1 A P_4 = B$
		\item $P_2 A P_3 = B$
		\item $P_2 A P_4 = B$
	\end{choices}
	
	
	\item 下列三维向量组不能构成 $R^3$ 的基的是 ( \ \ \ )
	\vspace{0.5em}
	\begin{choices}
		\item $\vec{\alpha}_1 = \begin{pmatrix} 1 \\ 0 \\ 0 \end{pmatrix}, \vec{\alpha}_2 = \begin{pmatrix} 1 \\ 2 \\ 0 \end{pmatrix}, \vec{\alpha}_3 = \begin{pmatrix} 1 \\ 2 \\ 3 \end{pmatrix}$
		\item $\vec{\alpha}_1 = \begin{pmatrix} 2 \\ 2 \\ 2 \end{pmatrix}, \vec{\alpha}_2 = \begin{pmatrix} 1 \\ 2 \\ 0 \end{pmatrix}, \vec{\alpha}_3 = \begin{pmatrix} 1 \\ 2 \\ 3 \end{pmatrix}$
		\item $\vec{\alpha}_1 = \begin{pmatrix} 2 \\ 5 \\ 4 \end{pmatrix}, \vec{\alpha}_2 = \begin{pmatrix} 3 \\ 6 \\ 9 \end{pmatrix}, \vec{\alpha}_3 = \begin{pmatrix} 1 \\ 2 \\ 3 \end{pmatrix}$
		\item $\vec{\alpha}_1 = \begin{pmatrix} 0 \\ 0 \\ 4 \end{pmatrix}, \vec{\alpha}_2 = \begin{pmatrix} 0 \\ 6 \\ 9 \end{pmatrix}, \vec{\alpha}_3 = \begin{pmatrix} 1 \\ 2 \\ 3 \end{pmatrix}$
	\end{choices}
	
	\item 设 $A, B, C$ 为同阶方阵, 下列矩阵的运算性质叙述正确的是 ( \ \ \ )
	\vspace{0.5em}
	\begin{choices}
		\item $(AB)^T = A^T B^T$
		\item $(A+B)^{-1} = A^{-1} + B^{-1}$
		\item $|AB| = |BA|$
		\item 若 $AB = AC$, 且 $A \neq O$, 则 $B = C$
	\end{choices}
	
	\item 行列式 $\begin{vmatrix} -1 & 0 & 2 \\ 3 & 4 & 5 \\ -5 & 8 & 7 \end{vmatrix}$ 中, $M_{11} + M_{12} + M_{13} = ( \ \ \ )$.
	\vspace{0.5em}
	\begin{choices}
		\item $78$
		\item $68$
		\item $84$
		\item $100$
	\end{choices}
	
	\item 在函数 $f(x) = \begin{vmatrix} 2x & 1 & -1 \\ -x & -x & x \\ 1 & x & 2 \end{vmatrix}$ 中, $x^3$ 的系数为 ( \ \ \ ).
	\vspace{0.5em}
	\begin{choices}
		\item $1$
		\item $-1$
		\item $2$
		\item $-2$
	\end{choices}
	
	\item 设 $A, B$ 均为 $n$ 阶矩阵, 则必有 ( \ \ \ )
	\vspace{0.5em}
	\begin{choices}
		\item $|A+B| = |A|+|B|$
		\item $|AB| = |BA|$
		\item $AB = BA$
		\item $(AB)^{-1} = A^{-1}B^{-1}$
	\end{choices}
	
	\item 下列命题正确的是 ( \ \ \ )
	\vspace{0.5em}
	\begin{choices}
		\item 一般地, $AB \neq BA$, 即使 $AB$ 和 $BA$ 都有意义.
		\item 若 $AB = O$ (零矩阵), 则 $A=O$ 或 $B=O$.
		\item 若 $AB=AC$, 且 $A \neq O$, 则 $B = C$.
		\item 若 $A(B-C)=O$, 且 $A \neq O$, 则 $B = C$.
	\end{choices}
	
	\item 已知向量组 $\vec{\alpha}_1 = \begin{pmatrix} 1 \\ 1 \\ 1 \end{pmatrix}$, $\vec{\alpha}_2 = \begin{pmatrix} 1 \\ 2 \\ 3 \end{pmatrix}$, $\vec{\alpha}_3 = \begin{pmatrix} 2 \\ 3 \\ c \end{pmatrix}$ 线性相关, 则 $c = ( \ \ \ )$
	\vspace{0.5em}
	\begin{choices}
		\item $2$
		\item $4$
		\item $6$
		\item $3$
	\end{choices}
	
	\item 设方阵 $A$ 满足 $|A-2E| = |2A+3E| = |3A-2E| = 0$, 则下列结论中正确的是 ( \ \ \  )
	\vspace{0.5em}
	\begin{choices}
		\item $|A| = -1$
		\item 方阵 $A$ 的特征值分别为 $1, 2, -\frac{3}{2}$
		\item 方阵 $A$ 可逆
		\item 方阵 $A$ 不可逆
	\end{choices}
	
	\item 向量组 $A: \vec{a}_1, \vec{a}_2, \dots, \vec{a}_s$ 线性相关的充要条件是 ( \ \ \ )
	\vspace{0.5em}
	\begin{choices}
		\item $A$ 中每个向量都可由其余向量线性表示
		\item $A$ 中至少有一个向量可由其余向量线性表示
		\item $A$ 中只有一个向量可由其余向量线性表示
		\item $A$ 中不包含零向量
	\end{choices}
	
	\item 设矩阵 $P = \begin{pmatrix} \frac{1}{\sqrt{2}} & -\frac{1}{\sqrt{2}} & 0 \\ \frac{1}{\sqrt{2}} & \frac{1}{\sqrt{2}} & 0 \\ 0 & 0 & 1 \end{pmatrix}$, $\vec{\alpha} = \begin{pmatrix} 1 \\ 2 \\ 3 \end{pmatrix}$, $\vec{\beta} = P\vec{\alpha}$, 则向量 $\vec{\beta}$ 的长度等于 ( \ \ \ )
	\vspace{0.5em}
	\begin{choices}
		\item $3$
		\item $\sqrt{6}$
		\item $14$
		\item $\sqrt{14}$
	\end{choices}
	
	\item 设 $A$ 为 $n$ 阶方阵, $\lambda_1, \lambda_2$ ( $\lambda_1 \neq \lambda_2$) 是 $A$ 的特征值, $\vec{\xi}_1, \vec{\xi}_2$ 是 $A$ 的分别属于特征值 $\lambda_1, \lambda_2$ 的特征向量. 则下列结论中错误的是 ( \ \ \ )
	\vspace{0.5em}
	\begin{choices}
		\item $k_1\vec{\xi}_1$ ($k_1 \neq 0$) 属于特征值 $\lambda_1$ 的特征向量
		\item $k_2\vec{\xi}_2$ ($k_2 \neq 0$) 属于特征值 $\lambda_2$ 的特征向量
		\item $\vec{\xi}_1, \vec{\xi}_2$ 线性相关
		\item $\vec{\xi}_1, \vec{\xi}_2$ 线性无关
	\end{choices}
	
	\item 分块对角矩阵 $A = \begin{pmatrix} 5 & 0 & 0 \\ 0 & 2 & 1 \\ 0 & 4 & 3 \end{pmatrix}$ 的逆矩阵为( \ \ \ )
	\vspace{0.5em}
	\begin{choices}
		\item $A = \begin{pmatrix} \frac{1}{5} & 0 & 0 \\ 0 & 2 & -1 \\ 0 & -4 & 3 \end{pmatrix}$
		\item $A = \begin{pmatrix} \frac{1}{5} & 0 & 0 \\ 0 & 1 & -\frac{1}{2} \\ 0 & -2 & \frac{3}{2} \end{pmatrix}$
		\item $A = \begin{pmatrix} \frac{1}{5} & 0 & 0 \\ 0 & -\frac{1}{2} & \frac{1}{2} \\ 0 & 2 & -1 \end{pmatrix}$
		\item $A = \begin{pmatrix} \frac{1}{5} & 0 & 0 \\ 0 & \frac{3}{2} & -\frac{1}{2} \\ 0 & -2 & 1 \end{pmatrix}$
	\end{choices}
	
	\item 已知四阶阵 $A$ 有特征值 $0, 1, 2, 3$, 则方程组 $A\vec{X}=\vec{0}$ 的基础解系所含向量个数为 ( \ \ \ )
	\vspace{0.5em}
	\begin{choices}
		\item $1$
		\item $2$
		\item $3$
		\item $4$
	\end{choices}
	
	\item 设矩阵 $A = \begin{pmatrix} 2 & 3 \\ 1 & 4 \end{pmatrix}$, 则 $A$ 的伴随矩阵 $A^*$ 为( \ \ \ )
	\vspace{0.5em}
	\begin{choices}
		\item $\begin{pmatrix} 4 & -3 \\ -1 & 2 \end{pmatrix}$
		\item $\begin{pmatrix} 4 & -1 \\ -3 & 2 \end{pmatrix}$
		\item $\begin{pmatrix} -2 & -3 \\ -1 & -4 \end{pmatrix}$
		\item $\begin{pmatrix} 2 & -3 \\ -1 & 4 \end{pmatrix}$
	\end{choices}
	
	\item 下列关于矩阵的性质错误的是( \ \ \ )
	\vspace{0.5em}
	\begin{choices}
		\item $R(A) = R(A^T)$
		\item 若 $A \sim* B$, 则 $R(A) = R(B)$
		\item $R(kA) = R(A)$ (k是实数)
		\item $\max(R(A), R(B)) \leq R(A, B) \leq R(A) + R(B)$
	\end{choices}
	
	\item 设 $A = \begin{pmatrix} 3 & -1 & 1 \\ 2 & 0 & -1 \\ 1 & -1 & 2 \end{pmatrix}$, 则 $A$ 的属于特征值 $\lambda=2$ 的一个特征向量是 ( \ \ \ )
	\vspace{0.5em}
	\begin{choices}
		\item $\begin{pmatrix} 1 \\ 0 \\ 1 \end{pmatrix}$
		\item $\begin{pmatrix} 1 \\ 0 \\ -1 \end{pmatrix}$
		\item $\begin{pmatrix} 1 \\ 1 \\ 0 \end{pmatrix}$
		\item $\begin{pmatrix} 0 \\ 1 \\ 1 \end{pmatrix}$
	\end{choices}
	
	\item 设三阶矩阵 $A = \begin{pmatrix} 1 & -\frac{1}{2} & \frac{1}{3} \\ -\frac{1}{2} & 1 & \frac{1}{2} \\ \frac{1}{3} & \frac{1}{2} & -1 \end{pmatrix}$ 和 $B = \begin{pmatrix} \frac{1}{9} & -\frac{8}{9} & -\frac{4}{9} \\ -\frac{8}{9} & \frac{1}{9} & -\frac{4}{9} \\ -\frac{4}{9} & -\frac{4}{9} & \frac{7}{9} \end{pmatrix}$, 下列选项正确的是 ( \ \ \ )
	\vspace{0.5em}
	\begin{choices}
		\item 这两个矩阵都是正交矩阵.
		\item 这两个矩阵都不是正交矩阵.
		\item 第一个矩阵是正交矩阵, 第二个矩阵不是正交矩阵.
		\item 第一个矩阵不是正交矩阵, 第二个矩阵是正交矩阵.
	\end{choices}
	
	\item 设有向量组 $\vec{\alpha}_1=(1,-1,2,4)^T$, $\vec{\alpha}_2=(0,3,1,2)^T$, $\vec{\alpha}_3=(3,0,7,14)^T$, $\vec{\alpha}_4=(1,-2,2,0)^T$, $\vec{\alpha}_5=(2,1,5,10)^T$, 则该向量组的最大线性无关组是 ( \ \ \ )
	\vspace{0.5em}
	\begin{choices}
		\item $\vec{\alpha}_1, \vec{\alpha}_2, \vec{\alpha}_3$
		\item $\vec{\alpha}_1, \vec{\alpha}_2, \vec{\alpha}_4$
		\item $\vec{\alpha}_1, \vec{\alpha}_2, \vec{\alpha}_5$
		\item $\vec{\alpha}_1, \vec{\alpha}_2, \vec{\alpha}_4, \vec{\alpha}_5$
	\end{choices}
	
	\item 设 $A$ 为3阶方阵, 已知 $|A|=-3$, $A^*$ 为 $A$ 的伴随矩阵,若交换A的第一行和第二行得到B, 则 $|BA^*| = ( \ \ \ )$
	\vspace{0.5em}
	\begin{choices}
		\item $9$
		\item $-9$
		\item $27$
		\item $-27$
	\end{choices}
	
	\item 四阶行列式 $D = \begin{vmatrix} -9 & 1 & 1 & 1 \\ 2 & -8 & 2 & 2 \\ 3 & 3 & -7 & 3 \\ 4 & 4 & 4 & -6 \end{vmatrix} = ( \ \ \ )$
	\vspace{0.5em}
	\begin{choices}
		\item $0$
		\item $1$
		\item $2$
		\item $3$
	\end{choices}
	
	
	\item 设 $A$ 为 $n$ 阶可逆阵, $\lambda$ 是 $A$ 的特征值, 则 $A^*$ 的特征值之一是( \ \ \ )
	\vspace{0.5em}
	\begin{choices}
		\item $\lambda^{-1} |A|^n$
		\item $\lambda^{-1} |A|$
		\item $\lambda |A|^n$
		\item $\lambda^{-1} |A|$
	\end{choices}
	
	\item 设 $A$ 是一个 $n(n \geq 3)$ 阶方阵, 下列陈述中正确的是 ( \ \ \ )
	\vspace{0.5em}
	\begin{choices}
		\item 若存在数 $\lambda$ 和向量 $\vec{\alpha}$ 使 $A\vec{\alpha} = \lambda\vec{\alpha}$, 则 $\vec{\alpha}$ 是 $A$ 的属于特征值 $\lambda$ 的特征向量;
		\item 若存在数 $\lambda$ 和非零向量 $\vec{\alpha}$ 使 $(A-\lambda E)\vec{\alpha} = \vec{0}$, 则 $\lambda$ 是 $A$ 的特征值;
		\item 矩阵 $A$ 的 $2$ 个不同特征值可以有同一个特征向量;
		\item 若 $\vec{\xi}$ 是矩阵 $A$ 的属于特征值 $\lambda$ 的特征向量, 则 $k\vec{\xi}$ 也是 $A$ 的特征向量.
	\end{choices}
	
	
	\item 已知 $\vec{\alpha} = (1, 2, 4, -2)^T$, $\vec{\beta} = (2, -4, -1, 2)^T$, 则 $[\vec{\alpha}, \vec{\beta}] = ( \ \ \ )$
	\vspace{0.5em}
	\begin{choices}
		\item $-14$
		\item $-12$
		\item $-10$
		\item $-16$
	\end{choices}
	
	\item 设 $A$ 是 $n$ 阶方阵, 且方程组 $A\vec{x}=\vec{0}$ 有非零解, 则矩阵 $A$ 必有一个特征值为 ( \ \ \ )
	\vspace{0.5em}
	\begin{choices}
		\item $-1$.
		\item $1$.
		\item $0$.
		\item $2$.
	\end{choices}
	
	\item 设 $A = \begin{pmatrix} 1 \\ 2 \end{pmatrix}$, $B = \begin{pmatrix} 3 & 4 \end{pmatrix}$, 则 $BA = ( \ \ \ )$
	\vspace{0.5em}
	\begin{choices}
		\item $7$
		\item $10$
		\item $11$
		\item $12$
	\end{choices}
	
	\item 设 $A$ 为 $3$ 阶矩阵, 且 $|3A+2E| = 0$, 则 $A$ 必有一个特征值为 ( \ \ \ )
	\vspace{0.5em}
	\begin{choices}
		\item $-\frac{3}{2}$
		\item $\frac{2}{3}$
		\item $-\frac{2}{3}$
		\item $\frac{3}{2}$
	\end{choices}
	
	\item 已知向量 $\vec{x}=(1,3,-2)^T$与$\vec{y}=(2,t,-5)^T$正交,则 $t=( \ \ \ )$
	\vspace{0.5em}
	\begin{choices}
		\item $4$
		\item $-4$
		\item $3$
		\item $-3$
	\end{choices}
	
	\item 若解向量 $\vec{\xi}_1, \vec{\xi}_2$ 是 $5$ 元齐次线性方程组 $A\vec{x}=\vec{0}$ 的基础解系, 则该方程组系数矩阵 $A$ 的秩为 ( \ \ \ )
	\vspace{0.5em}
	\begin{choices}
		\item $2$
		\item $3$
		\item $4$
		\item $5$
	\end{choices}
	
	\item 已知向量组 $A: \vec{\alpha}_1 = \begin{pmatrix} 1 \\ 0 \\ 1 \end{pmatrix}, \vec{\alpha}_2 = \begin{pmatrix} 1 \\ 1 \\ 1 \end{pmatrix}, \vec{\alpha}_3 = \begin{pmatrix} 1 \\ 2 \\ 3 \end{pmatrix}$ 和向量 $\vec{\beta} = \begin{pmatrix} 2 \\ 5 \\ 4 \end{pmatrix}$, 则下列结论中正确的是 ( )
	\vspace{0.5em}
	\begin{choices}
		\item 向量 $\vec{\beta}$ 不能由向量组 $A$ 线性表示
		\item 向量 $\vec{\beta}$ 能由向量组 $A$ 线性表示, 且表示式唯一
		\item 向量 $\vec{\beta}$ 能由向量组 $A$ 线性表示, 且表示式不唯一
		\item 以上说法都不正确
	\end{choices}
	
	\item 向量 $\vec{\alpha}=(1,2,0,2)^T, \vec{\beta}=(2,0,-1,2)^T$ 的夹角为 ( \ \ \ )
	\vspace{0.5em}
	\begin{choices}
		\item $\frac{\pi}{2}$.
		\item $\frac{\pi}{3}$.
		\item $\arcsin \frac{2}{3}$
		\item $\arccos \frac{2}{3}$
	\end{choices}
	
	\item 设矩阵 $A=\begin{pmatrix} 0 & 0 & 1 \\ x & 1 & y \\ 1 & 0 & 0 \end{pmatrix}$ 能对角化, 则 $x$ 和 $y$ 应满足条件 ( \ \ \ )
	\vspace{0.5em}
	\begin{choices}
		\item $x+y=0$.
		\item $x \neq y$.
		\item $x=y$.
		\item $x+y=1$.
	\end{choices}
	
	\item 设向量组 $A: \vec{\alpha}_1, \vec{\alpha}_2, \dots, \vec{\alpha}_m$ 可由向量组 $B: \vec{\beta}_1, \vec{\beta}_2, \dots, \vec{\beta}_s$ 线性表示, 且 $R_A = r_1, R_B = r_2$, 则下列选项正确的是( \ \ \ )
	\vspace{0.5em}
	\begin{choices}
		\item $r_1 = r_2$.
		\item $r_1 \leq r_2$.
		\item $r_1 \geq r_2$.
		\item 无法判定.
	\end{choices}
	
	\item 设 $A$ 为 $n$ 阶方阵, 则 $|A^*| = ( \ \ \ )$
	\vspace{0.5em}
	\begin{choices}
		\item $|A|$
		\item $0$
		\item $1$
		\item $|A|^{n-1}$
	\end{choices}
	
	\item 以下关于矩阵的秩的论述错误的是 ( \ \ \ )
	\vspace{0.5em}
	\begin{choices}
		\item 矩阵的秩是矩阵的最高阶非零子式的阶数.
		\item 矩阵化成行阶梯形矩阵后非零行的行数就是该矩阵的秩.
		\item 两个同型矩阵秩相等, 这两个矩阵就相等.
		\item 初等变换不改变矩阵的秩.
	\end{choices}
	
	\item 设 $A = \begin{pmatrix} -2 & 0 & 0 \\ 2 & x & 2 \\ 3 & 1 & 1 \end{pmatrix}$, $B = \begin{pmatrix} -1 & 0 & 0 \\ 0 & -2 & 0 \\ 0 & 0 & y \end{pmatrix}$, 且 $A$ 与 $B$ 相似, 则 ( \ \ \ )
	\vspace{0.5em}
	\begin{choices}
		\item $x=0, y=2$.
		\item $x=0, y=-2$.
		\item $x=2, y=0$.
		\item $x=-2, y=0$.
	\end{choices}
	
	\item 设矩阵 $A_{m \times n}$ 的秩为 $R(A) = m < n$, $I_m$ 为 $m$ 阶单位矩阵, 下述结论中正确的是 ( \ \ \ )
	\vspace{0.5em}
	\begin{choices}
		\item $A$ 的任意 $m$ 个列向量必线性无关.
		\item $A$ 的任意一个 $m$ 阶子式不等于零.
		\item 若矩阵 $B$ 满足 $BA = O$, 则 $B = O$.
		\item $A$ 通过初等行变换, 必可以化为 $(I_m \quad O)$ 的形式.
	\end{choices}
	
	\item 设 $A$ 为 $m \times n$ 矩阵, $R(A) < n$, 则下列关于齐次线性方程组 $A\vec{x}=\vec{0}$ 的叙述正确的是 ( \ \ \ )
	\vspace{0.5em}
	\begin{choices}
		\item $A\vec{x}=\vec{0}$ 只有零解.
		\item $A\vec{x}=\vec{0}$ 的基础解系含 $R(A)$ 个向量.
		\item $A\vec{x}=\vec{0}$ 的基础解系含 $n-R(A)$ 个向量.
		\item $A\vec{x}=\vec{0}$ 没有解.
	\end{choices}
	
	\item 齐次线性方程组 $\begin{cases} x_1+x_2+x_3=0 \\ 2x_1-x_2+\lambda x_3=0 \\ x_1-x_2+3x_3=0 \end{cases}$ 有非零解的充要条件是 $\lambda$ 取( \ \ \ )
	\vspace{0.5em}
	\begin{choices}
		\item $2$
		\item $3$
		\item $4$
		\item $5$
	\end{choices}
	
	\item 非齐次线性方程组 $\begin{cases} x_1 - x_2 + x_3 = 1 \\ x_1 + \lambda x_2 + x_3 = 1 \\ 2x_1 + 2\lambda x_2 + (\lambda+4)x_3 = 3 \end{cases}$, 有唯一解, 则 $\lambda$ 满足( \ \ \ )
	\vspace{0.5em}
	\begin{choices}
		\item $\lambda \neq -1$ 且 $\lambda \neq -2$
		\item $\lambda \neq -1$ 但 $\lambda = -2$
		\item $\lambda = -1$ 但 $\lambda \neq -2$
		\item $\lambda = -1$ 且 $\lambda = -2$
	\end{choices}
	
	\item 设 $A$ 为 $n$ 阶可逆矩阵, $\lambda$ 是 $A$ 的特征值, 则 $A^*$ 的特征值之一是 ( \ \ \ )
	\vspace{0.5em}
	\begin{choices}
		\item $\lambda^{-1} |A|$
		\item $\lambda^{-1} |A|^n$
		\item $\lambda |A|$
		\item $\lambda |A|^n$
	\end{choices}
	
	\item 齐次线性方程组 $\begin{cases} 2x_1 + \lambda x_2 + x_3 = 0 \\ (\lambda-1)x_1 - x_2 + 2x_3 = 0 \\ 4x_1 + x_2 + 4x_3 = 0 \end{cases}$ 有唯一解, 则 $\lambda = ( \ \ \ )$
	\vspace{0.5em}
	\begin{choices}
		\item $\lambda = 1$.
		\item $\lambda = 9/4$.
		\item $\lambda = 9$.
		\item $\lambda=1$ 或 $\lambda=9/4$.
	\end{choices}
	
	\item 排列 $3472516$ 的逆序数为 ( \ \ \ )
	\vspace{0.5em}
	\begin{choices}
		\item $7$
		\item $8$
		\item $9$
		\item $10$
	\end{choices}
	
	
	\item 设齐次线性方程组 $\begin{cases} x+y+z=0 \\ ax+by+cz=0 \\ bcx+cay+abz=0 \end{cases}$, 存在非零解, 则参数 $a, b, c$ 之间的关系是 ( \ \ \ )
	\vspace{0.5em}
	\begin{choices}
		\item $a=b=c$.
		\item $a=b$ 或 $b=c$ 或 $a=c$.
		\item $a, b, c$ 互不相等.
		\item $a \neq b$ 或 $b \neq c$ 或 $a \neq c$.
	\end{choices}
	
	\item 设非齐次线性方程组 $A\vec{x}=\vec{b}$ 无解, 矩阵 $A = \begin{pmatrix} 1 & 2 & 1 \\ 2 & 3 & a+2 \\ 1 & a & -2 \end{pmatrix}$ 常数项列向量 $\vec{b} = \begin{pmatrix} 1 \\ 3 \\ 0 \end{pmatrix}$, 则 $a= ( \ \ \ )$
	\vspace{0.5em}
	\begin{choices}
		\item $-2$
		\item $-1$
		\item $1$
		\item $2$
	\end{choices}
	
	\item 已知 $\vec{\alpha}_1 = (1, 1, 1)^T$, $\vec{\alpha}_2 = (0, 1, 1)^T$, $\vec{\alpha}_3 = (0, 0, 1)^T$ 是向量空间 $R^3$ 的一组基, 则向量 $\vec{\alpha} = (2, 0, 1)^T$ 在这组基下的坐标为 ( \ \ \ )
	\vspace{0.5em}
	\begin{choices}
		\item ($2, -2, 1$)
		\item ($-2, 2, 1$)
		\item ($1, -1, 1$)
		\item ($-1,1,-2$)
	\end{choices}
	
	\item 设 $A, B$ 均为 $n$ 阶方阵, 则必有( \ \ \ )
	\vspace{0.5em}
	\begin{choices}
		\item $|A+B| = |A|+|B|$
		\item $AB=BA$
		\item $|AB| = |BA|$
		\item $|A|^2 = |B|^2$
	\end{choices}
	
	\item 设 $A$ 是 $4 \times 3$ 矩阵, 且 $R(A)=2$, 而 $B = \begin{pmatrix} 1 & 0 & 2 \\ 0 & 2 & 0 \\ -1 & 0 & 3 \end{pmatrix}$, 则 $R(AB)=( \ \ \ )$
	\vspace{0.5em}
	\begin{choices}
		\item $1$
		\item $2$
		\item $3$
		\item 无法确定
	\end{choices}
	
	\item 设 $A = \begin{pmatrix} 1 & 1 & 0 \\ 1 & 0 & 1 \\ 0 & 1 & 1 \end{pmatrix}$, 则 $A$ 的特征值是 ( \ \ \  )
	\vspace{0.5em}
	\begin{choices}
		\item $0, 1, 1$
		\item $1, 1, 2$
		\item $-1, 1, 2$
		\item $-1, 1, 1$
	\end{choices}
	
	\item 设 $\vec{\alpha} = (1, 0, -2)^T$, $\vec{\beta} = (-4, 2, 3)^T$. 若 $\vec{\gamma}$ 与 $\vec{\alpha}$ 正交且 $\vec{\beta} = \lambda\vec{\alpha} + \vec{\gamma}$, 则 $\lambda = ( \ \ \ )$
	\vspace{0.5em}
	\begin{choices}
		\item $1$.
		\item $2$.
		\item $-1$.
		\item $-2$.
	\end{choices}
	
	\item 向量组 $\vec{\alpha}_1=(0,0,0,1)^T$, $\vec{\alpha}_2=(1,1,0,1)^T$, $\vec{\alpha}_3=(2,1,3,1)^T$, $\vec{\alpha}_4=(1,1,1,1)^T$ 的秩是 ( \ \ \ )
	\vspace{0.5em}
	\begin{choices}
		\item $1$
		\item $2$
		\item $3$
		\item $4$
	\end{choices}
	
	
	\item 若 $n$ 阶方阵 $A$ 与某对角阵相似, 则 ( \ \ \ )
	\vspace{0.5em}
	\begin{choices}
		\item $R(A) = n$.
		\item $A$ 有 $n$ 个互不相同的特征值.
		\item $A$ 有 $n$ 个线性无关的特征向量.
		\item $A$ 必为对称矩阵.
	\end{choices}
	
	\item 在四阶行列式 $D$ 中, $a_{11} = a_{12} = a_{13} = a_{14} = c \neq 0$. 以 $A_{ij}$ 表示元素 $a_{ij}$ 的代数余子式, 则 $A_{21} + A_{22} + A_{23} + A_{24} = ( \ \ \ )$
	\vspace{0.5em}
	\begin{choices}
		\item $0$
		\item $c$
		\item $-c$
		\item 不确定
	\end{choices}
	
	\item 设 $A$ 为 $3$ 阶方阵, $|A|=\frac{1}{2}$, 则$|(2A)^{-1}-5A^*|=$( \ \ \ )
	\vspace{0.5em}
	\begin{choices}
		\item $16$
		\item $-16$
		\item $4$
		\item $-54$
	\end{choices}
	
	\item 设 $\lambda = 2$ 是非零异号的 $A$ 的一个特征值, 则矩阵 $A^{-2}$ 有一个特征值为 ( \ \ \ )
	\vspace{0.5em}
	\begin{choices}
		\item $\frac{4}{3}$.
		\item $\frac{3}{4}$.
		\item $\frac{1}{2}$.
		\item $\frac{1}{4}$.
	\end{choices}
	
	\item 向量组 $\vec{\alpha}_1, \vec{\alpha}_2, \dots, \vec{\alpha}_s$ 线性无关的充要条件是 ( \ \ \ )
	\vspace{0.5em}
	\begin{choices}
		\item 向量组 $\vec{\alpha}_1, \vec{\alpha}_2, \dots, \vec{\alpha}_s$ 不含零向量.
		\item 向量组 $\vec{\alpha}_1, \vec{\alpha}_2, \dots, \vec{\alpha}_s$ 中任意两个向量线性无关.
		\item 向量组 $\vec{\alpha}_1$ 不能由向量组 $\vec{\alpha}_2, \dots, \vec{\alpha}_s$ 线性表示.
		\item 任一不全为零的数 $k_1, k_2, \dots, k_s$ 都不能使 $k_1\vec{\alpha}_1 + k_2\vec{\alpha}_2 + \dots + k_s\vec{\alpha}_s = \vec{0}$.
	\end{choices}
	
	\item $n$ 元非齐次线性方程组 $A\vec{X}=\vec{b}$ 无解的充要条件为 ( \ \ \ )
	\vspace{0.5em}
	\begin{choices}
		\item $R(A) = R(A, \vec{b})$
		\item $R(A) < R(A, \vec{b})$
		\item $R(A) \geq R(A, \vec{b})$
		\item $A$ 可逆
	\end{choices}
	
	\item 设矩阵 $A_{m \times n}$ 的秩为 $R(A)=m<n$, $E_m$ 为 $m$ 阶单位矩阵, 下述结论中正确的是( \ \ \ )
	\vspace{0.5em}
	\begin{choices}
		\item $A$ 的任意 $m$ 个列向量线性无关.
		\item $A$ 的任意一个 $m$ 阶子式不等于零.
		\item 若矩阵 $B$ 满足 $BA = O$, 则 $B = O$.
		\item $A$ 通过初等行变换, 可化成 $(E_m, O)$ 的形式.
	\end{choices}
	
	\item 若向量组 $\vec{\alpha}, \vec{\beta}, \vec{\gamma}$ 线性无关, $\vec{\alpha}, \vec{\beta}, \vec{\delta}$ 线性相关, 则 ( \ \ \ )
	\vspace{0.5em}
	\begin{choices}
		\item $\vec{\delta}$ 必可由 $\vec{\alpha}, \vec{\beta}, \vec{\gamma}$ 线性表示
		\item $\vec{\delta}$ 必不可由 $\vec{\alpha}, \vec{\beta}$ 线性表示
		\item $\vec{\gamma}$ 必可由 $\vec{\alpha}, \vec{\beta}, \vec{\delta}$ 线性表示
		\item $\vec{\gamma}$ 必不可由 $\vec{\alpha}, \vec{\beta}, \vec{\delta}$ 线性表示
	\end{choices}
	
	\item 用初等矩阵 $E = \begin{pmatrix} 1 & 0 & 0 \\ 0 & 0 & 1 \\ 0 & 1 & 0 \end{pmatrix}$ 左乘矩阵 $A = \begin{pmatrix} 2 & 1 & 1 \\ 3 & 1 & 1 \\ 2 & 4 & 6 \end{pmatrix}$ 相当于对矩阵 $A$ 进行的初等变换是( \ \ \ )
	\vspace{0.5em}
	\begin{choices}
		\item $r_1 \leftrightarrow r_2$
		\item $r_2 \leftrightarrow r_3$
		\item $c_1 \leftrightarrow c_2$
		\item $c_2 \leftrightarrow c_3$
	\end{choices}
	
	\item 设 $A$ 和 $B$ 有共同的特征值, 且各自有$n$个线性无关的特征向量, 则 ( \ \ \ )
	\vspace{0.5em}
	\begin{choices}
		\item $A = B$.
		\item $A \neq B$, 但 $|A-B|=0$.
		\item $A \sim* B$.
		\item $A$ 与 $B$ 不一定相似, 但 $|A|=|B|$.
	\end{choices}
	
	\item 如下不正确的是( \ \ \ )
	\vspace{0.5em}
	\begin{choices}
		\item $\begin{pmatrix} 0 & 0 & 1 \\ 0 & 1 & 0 \\ 1 & 0 & 0 \end{pmatrix}^{-1} = \begin{pmatrix} 0 & 0 & 1 \\ 0 & 1 & 0 \\ 1 & 0 & 0 \end{pmatrix}$
		\item $\begin{pmatrix} 1 & 0 & 0 \\ 0 & 3 & 0 \\ 0 & 0 & 1 \end{pmatrix}^{-1} = \begin{pmatrix} 1 & 0 & 0 \\ 0 & \frac{1}{3} & 0 \\ 0 & 0 & 1 \end{pmatrix}$
		\item $\begin{pmatrix} 1 & 0 & 0 \\ 0 & 1 & 0 \\ 2 & 0 & 1 \end{pmatrix}^{-1} = \begin{pmatrix} 1 & 0 & 2 \\ 0 & 1 & 0 \\ 0 & 0 & 1 \end{pmatrix}$
		\item $\begin{pmatrix} 1 & 0 & 0 \\ 0 & 1 & 0 \\ 2 & 0 & 1 \end{pmatrix}^{-1} = \begin{pmatrix} 1 & 0 & 0 \\ 0 & 1 & 0 \\ -2 & 0 & 1 \end{pmatrix}$
	\end{choices}
	
	\item 设 $3$ 阶行列式 $D$ 的第二列元素分别为 $1, 2, 3$, 对应的代数余子式分别为 $4, 5, 6$, 则该行列式第二列元素的代数余子式之和为 ( \ \ \ )
	\vspace{0.5em}
	\begin{choices}
		\item $-4+5-6$
		\item $4-5+6$
		\item $1 \times 4 + 2 \times 5 + 3 \times 6$
		\item $-1 \times 4 + 2 \times 5 - 3 \times 6$
	\end{choices}
	
	\item 设 $\vec{\eta}_1, \vec{\eta}_2$ 是非齐次线性方程组 $A\vec{x}=\vec{b}$ 的两个不同的解, 则 ( \ \ \ )
	\vspace{0.5em}
	\begin{choices}
		\item $\vec{\eta}_1+\vec{\eta}_2$ 是 $A\vec{x}=\vec{b}$ 的解.
		\item $\vec{\eta}_1-\vec{\eta}_2$ 是 $A\vec{x}=\vec{b}$ 的解.
		\item $2\vec{\eta}_1-3\vec{\eta}_2$ 是 $A\vec{x}=\vec{b}$ 的解.
		\item $3\vec{\eta}_1-2\vec{\eta}_2$ 是 $A\vec{x}=\vec{b}$ 的解.
	\end{choices}
	
	\item 设 $A$ 是 $m \times n$ 矩阵,若$m<n$, 则 ( \ \ \ )
	\vspace{0.5em}
	\begin{choices}
		\item 方程组 $A\vec{x}=\vec{b}$ 必有无穷多解.
		\item 方程组 $A\vec{x}=\vec{b}$ 无解.
		\item 方程组 $A\vec{x}=\vec{0}$ 必有无穷多解.
		\item 方程组 $A\vec{x}=\vec{0}$ 只有零解.
	\end{choices}
	
	\item 设 $A, B$ 是 $n$ 阶方阵, 且 $A$ 与 $B$ 相似, 则 ( \ \ \ )
	\vspace{0.5em}
	\begin{choices}
		\item $\lambda E - A = \lambda E - B$.
		\item $A$ 与 $B$ 有相同的特征值与特征向量.
		\item $A$ 与 $B$ 都相似于一个对角矩阵.
		\item 对任意常数 $k$, $kE-A$ 与 $kE-B$ 相似.
	\end{choices}
	
	\item 行列式 $\begin{vmatrix} 0 & 0 & 0 & 1 \\ 0 & 0 & 2 & 0 \\ 0 & 3 & 0 & 0 \\ 4 & 0 & 0 & 0 \end{vmatrix} = ( \ \ \ )$
	\vspace{0.5em}
	\begin{choices}
		\item $24$.
		\item $-24$.
		\item $10$.
		\item $-10$.
	\end{choices}
	
	\item 若 $n$ 阶方阵 $A$ 与对角阵 $\Lambda = \texttt{diag}(-1, 1, -1)$ 相似, 则 $A^6 = ( \ \ \ )$
	\vspace{0.5em}
	\begin{choices}
		\item $A$
		\item $-E$
		\item $E$
		\item $6E$
	\end{choices}
	
	
	\item 设 $\vec{\alpha}_1, \vec{\alpha}_2$ 分别是 $n$ 阶方阵 $A$ 的属于特征值 $\lambda_1, \lambda_2$ 的特征向量. 若 $\lambda_1 \neq \lambda_2$, 则 $\vec{\alpha}_1$ 与 $\vec{\alpha}_2$ ( \ \ \ )
	\vspace{0.5em}
	\begin{choices}
		\item 线性相关.
		\item 线性无关.
		\item 相等.
		\item 正交.
	\end{choices}
	
	
	\item 已知 $A = \begin{pmatrix} 1 & -2 & 1 \\ 2 & -1 & 2 \\ -1 & -2 & 4 \end{pmatrix}$, $B = \begin{pmatrix} 2 & -1 & 1 \\ 3 & 1 & -4 \\ 1 & -2 & -1 \end{pmatrix}$. $C = AB = (c_{ij})$, 则 $c_{32} = ( \ \ \ )$
	\vspace{0.5em}
	\begin{choices}
		\item $-9$
		\item $-11$
		\item $4$
		\item $2$
	\end{choices}
	
	\item 已知 $\vec{\xi}_1, \vec{\xi}_2$ 是 $A_{m×n}\vec{x}=\vec{0}$ 的两个不同解向量, 且 $R(A)=n-1$,$k$是任意常数,则该方程组的通解为( \ \ \ )
	\vspace{0.5em}
	\begin{choices}
		\item $k(\vec{\xi}_1+\vec{\xi}_2)$.
		\item $k(\vec{\xi}_1-\vec{\xi}_2)$.
		\item $k\vec{\xi}_1$.
		\item 不知道.
	\end{choices}
	
	\item 设 $A$ 为 $n$ 阶方阵, $k$ 为非零常数, 则 $|kA| = ( \ \ \ )$
	\vspace{0.5em}
	\begin{choices}
		\item $k|A|$.
		\item $|k||A|$.
		\item $k^n|A|$.
		\item $|k|^n|A|$.
	\end{choices}
	
	\item 设 $A$ 与 $B$ 相似, 则( \ \ \ )
	\vspace{0.5em}
	\begin{choices}
		\item $A$ 与 $B$ 有相同的特征向量.
		\item $A$ 与 $B$ 都和同一对角矩阵相似.
		\item $A - \lambda E = B - \lambda E$.
		\item $|A|=|B|$.
	\end{choices}
	
	\item 设 $\vec{\alpha}_1 = (1, 1, -1)^T$, $\vec{\alpha}_2 = (-2, -1, 2)^T$, 向量 $\vec{\alpha} = (2, \lambda, \mu)^T$ 与 $\vec{\alpha}_1, \vec{\alpha}_2$ 都正交, 则 $\lambda = ( \ \ \ )$
	\vspace{0.5em}
	\begin{choices}
		\item $1$
		\item $2$
		\item $0$
		\item $3$
	\end{choices}
	
	\item 设 $A$ 是 $n$ 阶矩阵, $\lambda_1, \lambda_2 (\lambda_1 \neq \lambda_2)$ 是 $A$ 的特征值, $\vec{\xi}_1, \vec{\xi}_2$ 是 $A$ 的分别属于 $\lambda_1, \lambda_2$ 的特征向量, 则下列结论中正确的是 ( \ \ \ )
	\vspace{0.5em}
	\begin{choices}
		\item 向量 $\vec{\xi}_1, \vec{\xi}_2$ 的对应分量成比例.
		\item 若 $\lambda_1 = 0$, 则 $\vec{\xi}_1 = \vec{0}$.
		\item 若 $\lambda_3 = \lambda_1 + \lambda_2$ 也是 $A$ 的特征值, 则其对应的特征向量是 $\vec{\xi}_1 + \vec{\xi}_2$.
		\item 向量 $\vec{\xi}_1 + \vec{\xi}_2$ 不可能是 $A$ 的特征向量.
	\end{choices}
	
	\item 设 $A = \begin{pmatrix} 0 & 0 & 1 \\ 0 & 1 & 0 \\ 1 & 0 & 0 \end{pmatrix}$, 则 $A$ 的特征值是 ( \ \ \ )
	\vspace{0.5em}
	\begin{choices}
		\item $-1, 1, 1$.
		\item $0, 1, 1$.
		\item $-1, 1, 2$.
		\item $1, 1, 2$.
	\end{choices}
	
	\item 设 $A = (\vec{\alpha}_1, \vec{\alpha}_2, \vec{\alpha}_3, \vec{\alpha}_4, \vec{\alpha}_5) = \begin{pmatrix} 1 & 3 & 2 & 1 & 3 \\ -1 & 1 & 0 & 1 & -1 \\ 1 & 1 & 1 & 0 & 2 \\ -1 & 3 & 1 & 2 & 0 \end{pmatrix}$, 则不是 $A$ 的列向量组的极大无关组的选项是 ( \ \ \ )
	\vspace{0.5em}
	\begin{choices}
		\item $\vec{\alpha}_1, \vec{\alpha}_3$
		\item $\vec{\alpha}_2, \vec{\alpha}_3$
		\item $\vec{\alpha}_2, \vec{\alpha}_4$
		\item $\vec{\alpha}_1, \vec{\alpha}_3, \vec{\alpha}_4$
	\end{choices}
	
	\item 能用克拉默法则求解的线性方程组须满足( \ \ \  )
	\vspace{0.5em}
	\begin{choices}
		\item 只须方程个数和未知数个数相等.
		\item 方程个数要等于未知数个数且系数行列式不等于 $0$.
		\item 方程个数可以不等于未知数个数.
		\item 系数行列式可以不等于 $0$.
	\end{choices}
	
	\item 对齐次线性方程组 $A\vec{x}=\vec{0}$, 若方程个数小于未知数个数, 则 ( \ \ \ )
	\vspace{0.5em}
	\begin{choices}
		\item $A\vec{x}=\vec{0}$ 无解.
		\item $A\vec{x}=\vec{0}$ 有唯一解.
		\item $A\vec{x}=\vec{0}$ 有非零解.
		\item $A\vec{x}=\vec{0}$ 只有零解.
	\end{choices}
	
	\item 设 $A$ 为 $n$ 阶方阵, 则 $|A^*| = ( \ \ \ )$
	\vspace{0.5em}
	\begin{choices}
		\item $|A|$
		\item $1$
		\item $0$
		\item $|A|^{n-1}$
	\end{choices}
	
	\item 二阶行列式 $\begin{vmatrix} \cos\theta & -\sin\theta \\ \sin\theta & \cos\theta \end{vmatrix} = ( \ \ \ )$
	\vspace{0.5em}
	\begin{choices}
		\item $1$
		\item $\cos 2\theta$
		\item $-1$
		\item $0$
	\end{choices}
	
	\item 下列命题不正确的是( \ \ \ )
	\vspace{0.5em}
	\begin{choices}
		\item 矩阵的秩等于它的列向量组的秩, 也等于它的行向量组的秩.
		\item 等价的向量组具有相同的秩.
		\item 若矩阵 $A$ 与 $B$ 行等价, 则 $A$ 中任意$s$个列向量与 $B$ 中对应的列向量具有相同的线性关系.
		\item 若矩阵 $A$ 与 $B$ 列等价, 则 $A$ 中任意$s$个列向量与 $B$ 中对应的列向量具有相同的线性关系.
	\end{choices}
	
	
	\item 矩阵 $A = \begin{pmatrix} 1 & 2 & -2 \\ 4 & t & -3 \\ 3 & -1 & 1 \end{pmatrix}$ 是降秩矩阵, 则 $t = ( \ \ \ )$
	\vspace{0.5em}
	\begin{choices}
		\item $1$
		\item $2$
		\item $3$
		\item $4$
	\end{choices}
	
	\item 已知向量组 $\vec{\alpha}_1 = \begin{pmatrix} k \\ 1 \\ 1 \end{pmatrix}$, $\vec{\alpha}_2 = \begin{pmatrix} 0 \\ 2 \\ 3 \end{pmatrix}$, $\vec{\alpha}_3 = \begin{pmatrix} 1 \\ 2 \\ 1 \end{pmatrix}$, 且 $R_A = 2$, 则 $k = ( \ \ \ )$
	\vspace{0.5em}
	\begin{choices}
		\item $0$
		\item $1$
		\item $2$
		\item $\frac{1}{4}$
	\end{choices}
	
	\item 下列矩阵中, 不是初等矩阵的是( \ \ \ )
	\vspace{0.5em}
	\begin{choices}
		\item $A = \begin{pmatrix} 0 & 0 & 1 \\ 0 & 1 & 0 \\ 1 & 0 & 0 \end{pmatrix}$
		\item $A = \begin{pmatrix} 1 & 0 & 0 \\ 0 & 2 & 0 \\ 0 & 0 & 1 \end{pmatrix}$
		\item $A = \begin{pmatrix} 1 & 0 & 0 \\ 0 & 0 & 0 \\ 0 & 1 & 0 \end{pmatrix}$
		\item $A = \begin{pmatrix} 1 & 0 & -2 \\ 0 & 1 & 0 \\ 0 & 0 & 1 \end{pmatrix}$
	\end{choices}
	
	
	\item 设向量组 $\vec{\alpha}_1, \vec{\alpha}_2, \dots, \vec{\alpha}_m$ 的秩为 $r$, 则 ( \ \ \  )
	\vspace{0.5em}
	\begin{choices}
		\item $\vec{\alpha}_1, \vec{\alpha}_2, \dots, \vec{\alpha}_m$ 中至少有一个由 $r$ 个向量组成的部分组线性无关.
		\item $\vec{\alpha}_1, \vec{\alpha}_2, \dots, \vec{\alpha}_m$ 中存在由 $r+1$ 个向量组成的部分组线性无关.
		\item $\vec{\alpha}_1, \vec{\alpha}_2, \dots, \vec{\alpha}_m$ 中由 $r$ 个向量组成的部分组都线性无关.
		\item $\vec{\alpha}_1, \vec{\alpha}_2, \dots, \vec{\alpha}_m$ 中个数小于 $r$ 的任意部分组都线性无关.
	\end{choices}
	
	
	\item 行列式 $\begin{vmatrix} 1 & 1 & 1 & 1 \\ 2 & 3 & -1 & -4 \\ 4 & 9 & 1 & 16 \\ 8 & 27 & -1 & -64 \end{vmatrix}$ 的值为( \ \ \ )
	\vspace{0.5em}
	\begin{choices}
		\item $-504$
		\item $1512$
		\item $-1512$
		\item $504$
	\end{choices}
	
	\item 设矩阵 $A = (1, -1, 1)$, $B = (1, 2, 3)$, 则 $A^T B = ( \ \ \ )$
	\vspace{0.5em}
	\begin{choices}
		\item $2$
		\item $0$
		\item $\begin{pmatrix} 1 & 2 & 3 \\ -1 & -2 & -3 \\ 1 & 2 & 3 \end{pmatrix}$
		\item $\begin{pmatrix} 1 & -1 & 1 \\ 2 & -2 & 2 \\ 3 & -3 & 3 \end{pmatrix}$
	\end{choices}
	
	\item 若向量组 $\vec{\alpha}_1, \vec{\alpha}_2$ 线性相关, 则 ( \ \ \ )
	\vspace{0.5em}
	\begin{choices}
		\item 存在不为0的常数 $k$ 使得 $\vec{\alpha}_1 = k\vec{\alpha}_2$或$\vec{\alpha}_2 = k\vec{\alpha}_1$.
		\item $\vec{\alpha}_1$ 和 $\vec{\alpha}_2$ 中必有一个零向量.
		\item $\vec{\alpha}_1$ 和 $\vec{\alpha}_2$ 都是零向量.
		\item $\vec{\alpha}_1 = \vec{\alpha}_2$.
	\end{choices}
	
	\item 对于方程个数和未知数个数相等的齐次线性方程组, 下列结论正确的是( \ \ \ )
	\vspace{0.5em}
	\begin{choices}
		\item 方程组不会无解.
		\item 方程组可以无解.
		\item 系数行列式不等于 $0$, 方程组有非零解.
		\item 系数行列式等于 $0$, 方程组只有零解.
	\end{choices}
	
	\item 已知 $\vec{\xi}_1, \vec{\xi}_2$ 是 $4$ 元齐次线性方程组 $A\vec{x}=\vec{0}$ 的一个基础解系, 则下列说法错误的是( \ \ \ )
	\vspace{0.5em}
	\begin{choices}
		\item $A\vec{\xi}_1 = \vec{0}, A\vec{\xi}_2 = \vec{0}$
		\item $A(\vec{\xi}_1 + \vec{\xi}_2) = \vec{0}$
		\item $\vec{\xi}_1, \vec{\xi}_2$ 线性相关
		\item $R(A) = 2$
	\end{choices}
	
	\item 关于 $3$ 维向量, 下列叙述不正确的是( \ \ \ )
	\vspace{0.5em}
	\begin{choices}
		\item 一个 $3$ 维非零向量线性无关.
		\item 两个 $3$ 维非零向量线性相关, 则它们共线.
		\item 三个 $3$ 维非零向量线性相关, 则它们共面.
		\item 四个 $3$ 维非零向量可以线性无关.
	\end{choices}
	
	
	\item 设三阶行列式 $D_3 = |\vec{\alpha}, \vec{\beta}, \vec{\gamma}|$, $\vec{\alpha}, \vec{\beta}, \vec{\gamma}$ 分别表示 $D_3$ 的三个列向量, 则下列选项中仍等于 $D_3$ 的是 ( \ \ \  )
	\vspace{0.5em}
	\begin{choices}
		\item $|\vec{\gamma}, \vec{\beta}, \vec{\alpha}|$
		\item $|\vec{\alpha}+\vec{\beta}, \vec{\beta}+\vec{\gamma}, \vec{\gamma}+\vec{\alpha}|$
		\item $|-\vec{\alpha}, -\vec{\beta}, -\vec{\gamma}|$
		\item $|\vec{\alpha}, \vec{\alpha}+\vec{\beta}, \vec{\alpha}+\vec{\beta}+\vec{\gamma}|$
	\end{choices}
	
	\item 若 $n$ 阶方阵 $A$ 与某对角矩阵相似, 则 ( \ \ \ )
	\vspace{0.5em}
	\begin{choices}
		\item $R(A) = n$.
		\item $A$ 有两个不同的特征值.
		\item $A$ 有 $n$ 个线性无关的特征向量.
		\item $A$ 应为对称矩阵.
	\end{choices}
	
	\item 关于 $n$ 维向量的定义, 下列叙述不正确的是( \ \ \ )
	\vspace{0.5em}
	\begin{choices}
		\item $n$ 维向量是 $n$ 个数构成的有序数组.
		\item 同样的 $n$ 个数构成的 $n$ 维行向量和 $n$ 维列向量要看作两个不同的向量.
		\item $n$ 个数顺序不同构成的向量相同.
		\item $2, 3$ 维向量具有几何意义, $4$ 维以上的向量没有具体几何意义.
	\end{choices}
	
	\item 当方阵 $A$ 可逆时, 求 $A$ 的逆矩阵的方法不包括( \ \ \ )
	\vspace{0.5em}
	\begin{choices}
		\item 利用定义$AA^{-1}=E$.
		\item 利用公式 $A^{-1} = \frac{1}{|A|} A^*$.
		\item 对矩阵 $(A, E)$ 进行初等行变换可化为 $(E, A^{-1})$.
		\item 对矩阵 $(A, E)$ 进行初等列变换可化为 $(E, A^{-1})$.
	\end{choices}
	
	
	\item 已知 $4$ 阶方阵 $A = (\vec{\alpha}_1, \vec{\alpha}_2, \vec{\alpha}_3, \vec{\alpha}_4)$ 的秩为 $3$, 且 $\vec{\alpha}_1 = 2\vec{\alpha}_2 - \vec{\alpha}_3$; 若 $\vec{\beta} = \vec{\alpha}_1 + \vec{\alpha}_2 + \vec{\alpha}_4$, 则下列结论错误的是 ( \ \ \  )
	\vspace{0.5em}
	\begin{choices}
		\item $(1, -2, 1, 0)^T$ 为齐次线性方程组 $A\vec{x}=\vec{0}$ 的解.
		\item $(1, 1, 1, 1)^T$ 为非齐次线性方程组 $A\vec{x}=\vec{\beta}$ 的解.
		\item 线性方程组 $A\vec{x}=\vec{\beta}$ 的通解为 $\vec{x} = k(1, -2, 1, 0)^T + (1, 1, 1, 1)^T$ ($k$ 为任意常数).
		\item 线性方程组 $A\vec{x}=\vec{\beta}$ 的通解为 $\vec{x} = k(1, 1, 1, 1)^T + (1, -2, 1, 0)^T$ ($k$ 为任意常数).
	\end{choices}
	
	\item 设 $A = \begin{pmatrix} a_{11} & a_{12} & a_{13} \\ a_{21} & a_{22} & a_{23} \\ a_{31} & a_{32} & a_{33} \end{pmatrix}$ 为 $3$ 阶非奇异矩阵, 则齐次线性方程组 $A\vec{x}=\vec{0}$ 的解为 ( \ \ \ )
	\vspace{0.5em}
	\begin{choices}
		\item $(1, 1, 1)^T$
		\item $(0, 0, 0)^T$
		\item $(a_{11}, a_{12}, a_{13})^T$
		\item $(0, 1, 0)^T$
	\end{choices}
	
\end{enumerate}


\section{解答题}
\begin{enumerate}
	
	\item 计算行列式 $D =
	\begin{vmatrix}
		3 & 1 & -1 & 2 \\
		-5 & 1 & 3 & -4 \\
		2 & 0 & 1 & -1 \\
		1 & -5 & 3 & -3
	\end{vmatrix}$.
	
	\item 计算$n$阶行列式 $D_n =
	\begin{vmatrix}
		4 & 1 & 1 & \cdots & 1 \\
		1 & 4 & 1 & \cdots & 1 \\
		\vdots & \vdots & \vdots & \ddots & \vdots \\
		1 & 1 & 1 & \cdots & 4 \\
		1 & 1 & 1 & \cdots & 4
	\end{vmatrix}$.
	
	\item 计算四阶行列式 $D =
	\begin{vmatrix}
		1 & 2 & 3 & 4 \\
		1 & 0 & 1 & 2 \\
		3 & -1 & -1 & 0 \\
		1 & 2 & 0 & -5
	\end{vmatrix}$.
	
	\item 设 $D =
	\begin{vmatrix}
		3 & -5 & 2 & 1 \\
		1 & 1 & 0 & -5 \\
		-1 & 3 & 1 & 3 \\
		2 & -4 & -1 & -3
	\end{vmatrix}$
	,$D$ 的 $(i,j)$ 元的余子式和代数余子式依次记作 $M_{ij}$ 和 $A_{ij}$,求
	$A_{11} + A_{12} + A_{13} + A_{14}$ 及 $M_{11} + M_{21} + M_{31} + M_{41}$.
	
	
	
	\item 对于含参数$\lambda$的齐次线性方程组
	$
	\begin{cases}
		(1-\lambda)x_1 - 2x_2 + 4x_3 = 0, \\
		2x_1 + (3-\lambda)x_2 + x_3 = 0, \\
		x_1 + x_2 + (1-\lambda)x_3 = 0.
	\end{cases}
	$
	问$\lambda$取何值时该齐次线性方程组有非零解?
	
	
	
	\item 已知 $A =
	\begin{pmatrix}
		3 & 1 & 0 \\
		-1 & 2 & 1 \\
		3 & 4 & 2
	\end{pmatrix}$, $B =
	\begin{pmatrix}
		1 & 0 & 2 \\
		-1 & 1 & 1 \\
		2 & 1 & 1
	\end{pmatrix}$, 求满足方程 $3A - 2X = B$ 中的 $X$.
	
	
	
	\item 设矩阵 $A =
	\begin{pmatrix}
		1 & -3 & 2 \\
		-3 & 0 & 1 \\
		1 & 1 & -1
	\end{pmatrix}$, 求 $A^{-1}$.
	
	
	\item 解矩阵方程
	$
	\begin{pmatrix}
		0 & 1 & 0 \\
		1 & 0 & 0 \\
		0 & 0 & 1
	\end{pmatrix}
	X
	\begin{pmatrix}
		1 & 0 & 0 \\
		0 & 1 & -2 \\
		0 & 0 & 1
	\end{pmatrix}
	=
	\begin{pmatrix}
		1 & -4 & 3 \\
		2 & 0 & -1 \\
		1 & -2 & 0
	\end{pmatrix}.
	$
	
	
	\item 把矩阵
	$
	\begin{pmatrix}
		2 & 3 & 1 & -3 & -7 \\
		1 & 2 & 0 & -2 & -4 \\
		3 & -2 & 8 & 3 & 0 \\
		2 & -3 & 7 & 4 & 3
	\end{pmatrix}
	$
	化为行最简矩阵.
	
	
	
	\item 解线性方程组
	$
	\begin{cases}
		2x + 4y - z = 1, \\
		x - 3y + 2z = 3, \\
		3x + y + z = 5.
	\end{cases}
	$
	
	
	\item 已知 $A =
	\begin{pmatrix}
		1 & 1 & -1 \\
		0 & 1 & 1 \\
		0 & 0 & -1
	\end{pmatrix}$, 且 $A^2 - AB = E$, $E$ 是三阶单位矩阵, 求矩阵 $B$.
	
	
	\item 设矩阵 $A =
	\begin{pmatrix}
		1 & 1 & 1 & 1 & 1 \\
		3 & 2 & 1 & -3 & x \\
		0 & 1 & 2 & 6 & 3 \\
		5 & 4 & 3 & -1 & y
	\end{pmatrix}$
	的秩 $R(A)=2$, 求 $x$ 和 $y$ 的值.
	
	
	
	\item 设 $A =
	\begin{pmatrix}
		-1 & 1 & 0 \\
		0 & 0 & 2 \\
		0 & 0 & 2
	\end{pmatrix}$, $B =
	\begin{pmatrix}
		1 & 1 & 0 \\
		0 & 2 & 2 \\
		0 & 0 & 3
	\end{pmatrix}$, $X$ 是 $3$ 阶未知矩阵, 且
	$(E - B^{-1}A)^T B X = E$, 求 $X, X^{-1}$.
	
	
	\item 解非齐次线性方程组
	$
	\begin{cases}
		2x_1 + 7x_2 + 3x_3 + x_4 = 6, \\
		3x_1 + 5x_2 + 2x_3 + 2x_4 = 4, \\
		5x_1 + 12x_2 + 5x_3 + 3x_4 = 10.
	\end{cases}
	$,写出方程组的通解.
	
	
	\item 当 $\lambda$ 取何值时, 方程组
	$
	\begin{cases}
		\lambda x_1 + x_2 + x_3 = 1, \\
		x_1 + \lambda x_2 + x_3 = \lambda, \\
		x_1 + x_2 + \lambda x_3 = \lambda^2.
	\end{cases}
	$
	无解、有唯一解、有无穷多解？有无穷多解时求出通解.
	
	
	\item 设 $\vec{\alpha}_1 = \begin{pmatrix} 1 \\ 1 \\ 1 \end{pmatrix}$, $\vec{\alpha}_2 = \begin{pmatrix} 1 \\ 2 \\ 3 \end{pmatrix}$, $\vec{\alpha}_3 = \begin{pmatrix} 1 \\ 3 \\ t \end{pmatrix}$, 问为何值时, 向量组 $\vec{\alpha}_1, \vec{\alpha}_2, \vec{\alpha}_3$ 线性相关?
	$t$ 为何值时, 向量组 $\vec{\alpha}_1, \vec{\alpha}_2, \vec{\alpha}_3$ 线性无关?
	
	
	\item 已知 $\vec{\alpha}_1 = \begin{pmatrix} 1 \\ 1 \\ 0 \\ -5 \end{pmatrix}$, $\vec{\alpha}_2 = \begin{pmatrix} -1 \\ 3 \\ 1 \\ 3 \end{pmatrix}$, $\vec{\alpha}_3 = \begin{pmatrix} 2 \\ -4 \\ -1 \\ -3 \end{pmatrix}$, 判断 $\vec{\alpha}_1, \vec{\alpha}_2, \vec{\alpha}_3$ 是否线性相关?
	
	
	\item 设 $\vec{\alpha}_1 = \begin{pmatrix} 1 \\ 1 \\ 2 \\ 2 \end{pmatrix}$, $\vec{\alpha}_2 = \begin{pmatrix} 1 \\ 2 \\ 1 \\ 3 \end{pmatrix}$, $\vec{\alpha}_3 = \begin{pmatrix} 1 \\ -1 \\ 4 \\ 0 \end{pmatrix}$, $\vec{\beta} = \begin{pmatrix} 1 \\ 0 \\ 3 \\ 1 \end{pmatrix}$, 证明向量 $\vec{\beta}$ 能由向量组
	$\vec{\alpha}_1, \vec{\alpha}_2, \vec{\alpha}_3$ 线性表示, 并求出表达式.
	
	
	\item 设 $\vec{\alpha}_1 = \begin{pmatrix} 1 \\ -1 \\ 1 \\ -1 \end{pmatrix}$, $\vec{\alpha}_2 = \begin{pmatrix} 3 \\ 1 \\ 1 \\ 3 \end{pmatrix}$, $\vec{\beta}_1 = \begin{pmatrix} 2 \\ 0 \\ 1 \\ 1 \end{pmatrix}$, $\vec{\beta}_2 = \begin{pmatrix} 1 \\ 1 \\ 0 \\ 2 \end{pmatrix}$, $\vec{\beta}_3 = \begin{pmatrix} 3 \\ -1 \\ 2 \\ 0 \end{pmatrix}$, 证明向量组
	$\vec{\alpha}_1, \vec{\alpha}_2$ 与向量组 $\vec{\beta}_1, \vec{\beta}_2, \vec{\beta}_3$ 等价.
	
	
	\item 给定向量组 $\vec{\alpha}_1 = \begin{pmatrix} 2 \\ 1 \\ 4 \\ 3 \end{pmatrix}$, $\vec{\alpha}_2 = \begin{pmatrix} -1 \\ 1 \\ -6 \\ 6 \end{pmatrix}$, $\vec{\alpha}_3 = \begin{pmatrix} -1 \\ -2 \\ 2 \\ -9 \end{pmatrix}$, $\vec{\alpha}_4 = \begin{pmatrix} 1 \\ 1 \\ -2 \\ 7 \end{pmatrix}$, $\vec{\alpha}_5 = \begin{pmatrix} 2 \\ 4 \\ 4 \\ 9 \end{pmatrix}$, 求它
	的一个最大线性无关组, 并把不属于最大无关组的向量用最大无关组线性表示.
	
	
	\item 求齐次线性方程组
	$
	\begin{cases}
		x_1 + x_2 - x_3 - x_4 = 0, \\
		2x_1 - 5x_2 + 3x_3 + 2x_4 = 0, \\
		7x_1 - 7x_2 + 3x_3 + x_4 = 0.
	\end{cases}
	$
	的基础解系和通解.
	
	
	
	\item 求解非齐次线性方程组
	$
	\begin{cases}
		x_1 - x_2 - x_3 + x_4 = 0, \\
		x_1 - x_2 + x_3 - 3x_4 = 1, \\
		x_1 - x_2 - 2x_3 + 3x_4 = -\frac{1}{2}.
	\end{cases}
	$
	
	
	\item 试用施密特法把向量组 $\vec{\alpha}_1 = \begin{pmatrix} 1 \\ 1 \\ 1 \end{pmatrix}$、$\vec{\alpha}_2 = \begin{pmatrix} 1 \\ 2 \\ 3 \end{pmatrix}$、$\vec{\alpha}_3 = \begin{pmatrix} 2 \\ 1 \\ 3 \end{pmatrix}$ 正交化.
	
	
	\item 已知向量 $\vec{\alpha}_1 = \begin{pmatrix} 1 \\ 1 \\ 1 \end{pmatrix}$ 与 $\vec{\alpha}_2 = \begin{pmatrix} 1 \\ -2 \\ 1 \end{pmatrix}$ 正交, 求一个向量 $\vec{\alpha}_3$, 使 $\vec{\alpha}_1$、$\vec{\alpha}_2$、$\vec{\alpha}_3$ 两两正交.
	
	
	
	\item 已知四元非齐次线性方程组 $A\vec{x} = \vec{b}$ 的三个解是 $\vec{\eta}_1$、$\vec{\eta}_2$、$\vec{\eta}_3$, 且
	$\vec{\eta}_1 = (1, 2, 3, 4)^T$, $\vec{\eta}_2 + \vec{\eta}_3 = (3, 5, 7, 9)^T$, $R(A)=3$, 求方程组的通解.
	
	
	\item 设方阵 $A =
	\begin{pmatrix}
		2 & -1 & 2 \\
		5 & a & 3 \\
		-1 & b & -2
	\end{pmatrix}$ 有一个特征向量 $\vec{\alpha} = \begin{pmatrix} 1 \\ 1 \\ -1 \end{pmatrix}$, 求 $a, b$ 的值, 并求出对应于 $\vec{\alpha}$ 的特
	征值.
	
	
	\item 设方阵 $A =
	\begin{pmatrix}
		2 & 1 & 1 \\
		2 & 3 & 2 \\
		3 & 3 & 4
	\end{pmatrix}$, 求方阵 $A$ 的特征值和特征向量.
	
	
	
	\item 设三阶矩阵 $A$ 的特征值是 $-1$, $2$, $5$, 矩阵 $B = 3A - A^2$, 求$B$ 的特征值、$|B|$和$|A-3E|$.
	
	
	\item 设方阵 $A =
	\begin{pmatrix}
		-2 & 1 & 1 \\
		0 & 2 & 0 \\
		-4 & 1 & 3
	\end{pmatrix}$, 求 $g(A) = A^3 - 2A + E$ 的全部特征值和特征向量.
	
	
	\item 设方阵 $A =
	\begin{pmatrix}
		0 & 0 & 1 \\
		x & 1 & y \\
		1 & 0 & 0
	\end{pmatrix}$ 能够相似对角化, 求 $x, y$ 满足的条件.
	
	
	
	\item 已知矩阵 $A =
	\begin{pmatrix}
		1 & -1 & 1 \\
		2 & 4 & -2 \\
		-3 & -3 & x
	\end{pmatrix}$ 与 $B =
	\begin{pmatrix}
		2 & 0 & 0 \\
		0 & 2 & 0 \\
		0 & 0 & y
	\end{pmatrix}$ 相似, 求常数 $x, y$ 及可逆矩阵 $P$, 使得
	$P^{-1}AP = B$.
	
	
	\item 求一个正交的相似变换矩阵, 将对称矩阵 $A =
	\begin{pmatrix}
		2 & 2 & -2 \\
		2 & 5 & -4 \\
		-2 & -4 & 5
	\end{pmatrix}$ 化为对角阵.
	
	
	\item 设 $3$ 阶方阵 $A$ 的特征值为 $\lambda_1 = 2$, $\lambda_2 = -2$, $\lambda_3 = 1$, 对应的特征向量依次为
	$\vec{p}_1 = \begin{pmatrix} 0 \\ 1 \\ 1 \end{pmatrix}$, $\vec{p}_2 = \begin{pmatrix} 1 \\ 1 \\ 1 \end{pmatrix}$, $\vec{p}_3 = \begin{pmatrix} 1 \\ 1 \\ 0 \end{pmatrix}$, 求 $A$.
	
	
	
	\item 设 $A =
	\begin{pmatrix}
		3 & -2 \\
		-2 & 3
	\end{pmatrix}$, 求 $\varphi(A) = A^{10} - 5A^9$.
	
	
	\item 求矩阵 $A =
	\begin{pmatrix}
		2 & 0 & 1 \\
		3 & 1 & 3 \\
		4 & 0 & 5
	\end{pmatrix}$ 的全部特征值和特征向量, 判定矩阵能不能与对角阵相似, 说明理由.
	
	
	
	\item 设向量组 $(2, 1, 1, 1)$, $(2, 1, a, a)$, $(3, 2, 1, a)$, $(4, 3, 2, 1)$ 线性相关,且 $a \ne 1$,求 $a$.
	
	\item 已知方程组 
	$\begin{cases} 
		x_1 + 2x_2 + 3x_3 = 0 \\ 
		2x_1 + 3x_2 + 5x_3 = 0 \\
		x_1 + x_2 + ax_3 = 0 
	\end{cases}$ 
	与 
	$\begin{cases} 
		x_1 + bx_2 + cx_3 = 0 \\ 
		2x_1 + b^2x_2 + (c+1)x_3 = 0 
	\end{cases}$ 
	同解,求 $a, b, c$.
	
	\item 设 $A = \begin{pmatrix} 2 & 1 \\ -1 & 2 \end{pmatrix}$.若 $BA = B + 2E$,求 $|B|$.
	
	\item 设 $\alpha_1, \alpha_2, \dots, \alpha_s$ 为 $n$ 维列向量,$A$ 是 $m \times n$ 矩阵,判断当 $\alpha_1, \alpha_2, \dots, \alpha_s$ 线性相关或无关时,$A\alpha_1, A\alpha_2, \dots, A\alpha_s$ 的相关性.
	
	\item 设 $A$ 为 3 阶矩阵,将 $A$ 的第 2 行加到第 1 行得 $B$,再将 $B$ 的第 1 列的 $-1$ 倍加到第 2 列得 $C$,记 $P = \begin{pmatrix} 1 & 1 & 0 \\ 0 & 1 & 0 \\ 0 & 0 & 1 \end{pmatrix}$,试用 $A, P$ 表示 $C$.
	
	\item 设 4 维向量组 $\alpha_1 = (1+a, 1, 1, 1)$, $\alpha_2 = (2, 2+a, 2, 2)$, $\alpha_3 = (3, 3, 3+a, 3)$, $\alpha_4 = (4, 4, 4, 4+a)$.当 $a$ 为何值时,向量组线性无关？求其极大无关组,并用极大无关组表示其他向量.
	
	\item 设 $A = \begin{pmatrix} 0 & 1 & 0 & 0 \\ 0 & 0 & 1 & 0 \\ 0 & 0 & 0 & 1 \\ 0 & 0 & 0 & 0 \end{pmatrix}$,求 $\text{rank}(A^3)$.
	
	\item 设方程组 $\begin{cases} x_1 + x_2 + x_3 = 0 \\ x_1 + 2x_2 + ax_3 = 0 \\ x_1 + 4x_2 + a^2 x_3 = 0 \end{cases}$ 与 $ x_1 + 2x_2 + x_3 = a-1 $ 有公共解,求 $a$ 及所有公共解.
	
	\item $n$ 阶非零矩阵 $A$,若 $A^3 = 0$,判断 $E-A$ 与 $E+A$ 是否可逆.
	
	\item 设 $Ax = b$,其中系数矩阵 $A$ 为:
	\[ A = \begin{pmatrix} 
		2a & 1 &  &  &  \\ 
		a^2 & 2a & 1 &  &  \\ 
		& a^2 & 2a & 1 &   \\ 
		& & \ddots & \ddots & \ddots \\ 
		&   &  & a^2 & 2a & 1 \\
		& & & & a^2 & 2a
	\end{pmatrix}_{n \times n}, \quad x = \begin{pmatrix} x_1 \\ x_2 \\ \vdots \\ x_n \end{pmatrix}, \quad b = \begin{pmatrix} 1 \\ 0 \\ \vdots \\ 0 \end{pmatrix} \]
	\begin{enumerate}
		\item 证明:$|A| = (n+1)a^n$.
		\item 当 $a$ 为何值时,方程有唯一解,并求 $x$.
		\item 当 $a$ 为何值时,方程有无穷多解,并求通解.
	\end{enumerate}
	
	\item 设 $A, B$ 均为 2 阶矩阵,$A^*, B^*$ 分别为 $A, B$ 的伴随矩阵,若 $|A|=2, |B|=3$,求 $\begin{pmatrix} 0 & A \\ B & 0 \end{pmatrix}$ 的伴随矩阵.
	
	\item 设 $A, P$ 均为 3 阶矩阵,且 $P^{T}AP = \begin{pmatrix}
		1 \\
		& 1 \\
		& & 2
	\end{pmatrix}$.若 $P = (\alpha_1, \alpha_2, \alpha_3)$,$Q = (\alpha_1 + \alpha_2, \alpha_2, \alpha_3)$,求 $Q^T A Q$.
	
	\item 设 $A = \begin{pmatrix} 1 & -1 & -1 \\ -1 & 1 & 1 \\ 0 & -4 & -2 \end{pmatrix}$, $\alpha = \begin{pmatrix} -1 \\ 1 \\ -2 \end{pmatrix}$:
	\begin{enumerate}
		\item 求满足 $A\beta = \alpha, A^2 \eta = \alpha$ 的所有向量 $\beta, \eta$.
		\item 判断 $\alpha, \beta, \eta$ 的线性相关性.
	\end{enumerate}
	
	\item 设 3 阶矩阵 $A, B$ 满足 $|A|=3, |B|=2, |A^{-1}+B|=2$,求 $|A+B^{-1}|$.
	
	\item 设 $A = \begin{pmatrix} \lambda & 1 & 1 \\ 0 & \lambda-1 & 0 \\ 1 & 1 & \lambda \end{pmatrix}, b = \begin{pmatrix} a \\ 1 \\ 1 \end{pmatrix}$,已知 $Ax=b$ 有2个不同解,求 $\lambda, a$ 及通解.
	
	\item 设 $A$ 为 3 阶矩阵,将 $A$ 的第 2 列加到第 1 列得 $B$,再将 $B$ 的第 2 列与第 3 列交换得 $E$.记 $P_1 = \begin{pmatrix} 1 & 0 & 0 \\ 1 & 1 & 0 \\ 0 & 0 & 1 \end{pmatrix}, P_2 = \begin{pmatrix} 1 & 0 & 0 \\ 0 & 0 & 1 \\ 0 & 1 & 0 \end{pmatrix}$,试用 $P_1, P_2$ 表示 $A$.
	
	\item 设向量组 $\alpha_1=(1,0,1), \alpha_2=(0,1,1), \alpha_3=(1,3,5)$ 不能由向量组 $\beta_1=(1,1,1), \beta_2=(1,2,3), \beta_3=(3,4,a)$ 线性表示,求 $a$,并将 $\beta_1, \beta_2, \beta_3$ 用 $\alpha_1, \alpha_2, \alpha_3$ 表示.
	
	\item 设 $\alpha_1=(0,0,c_1), \alpha_2=(0,1,c_2), \alpha_3=(1,-1,c_3), \alpha_4=(-1,1,c_4)$,找出其中的线性相关组.
	
	\item 设 $A$ 为 3 阶矩阵,$P$ 为 3 阶可逆阵,且 $P^{-1}AP = \text{diag}(1, 1, 2)$.若 $P=(\alpha_1, \alpha_2, \alpha_3), Q=(\alpha_1+\alpha_2, \alpha_2, \alpha_3)$,求 $Q^{-1}AQ$.
	
	\item 设 $A$ 为 3 阶矩阵,$|A|=3$,$A^*$ 为 $A$ 的伴随矩阵.若交换 $A$ 的第 1 行与第 2 行得矩阵 $B$,求 $|BA^*|$.
	
	\item 设 $A = \begin{pmatrix} 1 & a & 0 & 0 \\ 0 & 1 & a & 0 \\ 0 & 0 & 1 & a \\ a & 0 & 0 & 1 \end{pmatrix}, b = \begin{pmatrix} 1 \\ -1 \\ 0 \\ 0 \end{pmatrix}$.求 $|A|$,及当 $Ax=b$ 有无穷多解时 $a$ 的值,并求通解.
	
	\item 设 $A=(a_{ij})$ 是 3 阶非零矩阵,$A_{ij}$ 为 $a_{ij}$ 的代数余子式.若 $a_{ij} + A_{ij} = 0$,求 $|A|$.
	
	\item 设 $A = \begin{pmatrix} 1 & a \\ 1 & 0 \end{pmatrix}, B = \begin{pmatrix} 0 & 1 \\ 1 & b \end{pmatrix}$.当 $a, b$ 为何值时,存在矩阵 $C$ 使得 $AC - CA = B$,求所有的 $C$.
	
	\item 计算行列式:$D = \begin{vmatrix} 0 & a & b & 0 \\ a & 0 & 0 & b \\ 0 & c & d & 0 \\ c & 0 & 0 & d \end{vmatrix}$.
	
	\item 设 $\alpha_1, \alpha_2, \alpha_3$ 是 3 维向量,对于任意常数 $l, k$,向量 $\alpha_1 + k\alpha_3, \alpha_2 + l\alpha_3$ 线性无关,是 $\alpha_1, \alpha_2, \alpha_3$ 线性无关的什么条件？
	
	\item 设 $A = \begin{pmatrix} 1 & -2 & 3 & -4 \\ 0 & 1 & -1 & 1 \\ 1 & 2 & 0 & -3 \end{pmatrix}$.求 $Ax=0$ 的基础解系,并求满足 $AB=E$ 的所有 $B$.
	
	\item 求 $\begin{pmatrix} 1 & 1 & 1 \\ 1 & 2 & a \\ 1 & 4 & a^2 \end{pmatrix}x = \begin{pmatrix} 1 \\ d \\ d^2 \end{pmatrix}$ 有无穷多解的充要条件.
	
	\item 设 $A = \begin{pmatrix} a & 1 & 0 \\ 1 & a & -1 \\ 0 & 1 & a \end{pmatrix}$,且 $A^3 = 0$,求 $a$ 及满足 $X - XA^2 - AX + AXA^2 = E$ 的 $X$.
	
	\item 计算行列式:$D = \begin{vmatrix} \lambda & -1 & 0 & 0 \\ 0 & \lambda & -1 & 0 \\ 0 & 0 & \lambda & -1 \\ 4 & 3 & 2 & \lambda + 1 \end{vmatrix}$.
	
	\item 设 $A = \begin{pmatrix} 1 & 1 & 1-a \\ 1 & 0 & a \\ a+1 & 1 & a+1 \end{pmatrix}, b = \begin{pmatrix} 0 \\ 1 \\ 2a-2 \end{pmatrix}$,且 $Ax=b$ 无解,求 $a$ 及 $A^T Ax = A^T b$ 的通解.
	
	\item 设 $A = \begin{pmatrix} 0 & -1 & 1 \\ 2 & -3 & 0 \\ 0 & 0 & 0 \end{pmatrix}$,求 $A^{99}$.若 $B = (\alpha_1, \alpha_2, \alpha_3)$ 且 $B^2 = BA$,记 $B^{100} = (\beta_1, \beta_2, \beta_3)$,将 $\beta_1, \beta_2, \beta_3$ 表示为 $\alpha_1, \alpha_2, \alpha_3$ 的线性组合.
	
	\item 设 $\alpha$ 是 $n$ 维单位列向量,判断 $E - \alpha\alpha^T, E + \alpha\alpha^T, E + 2\alpha\alpha^T, E - 2\alpha\alpha^T$ 是否可逆.
	
	\item 设 $A = \begin{pmatrix} 1 & 0 & 1 \\ 1 & 1 & 2 \\ 0 & 1 & 1 \end{pmatrix}$,$\alpha_1, \alpha_2, \alpha_3$ 为线性无关的 3 维列向量,求 $A\alpha_1, A\alpha_2, A\alpha_3$ 的秩.
	
	\item 设 $A$ 为 3 阶矩阵,$\alpha_1, \alpha_2, \alpha_3$ 线性无关.若 $A\alpha_1 = \alpha_1 + \alpha_2, A\alpha_2 = \alpha_2 + \alpha_3, A\alpha_3 = \alpha_1 + \alpha_3$,求 $|A|$.
	
	\item 设 $A = \begin{pmatrix} 1 & 2 & a \\ 1 & 3 & 0 \\ 2 & 7 & -a \end{pmatrix}$ 经初等列变换化为 $B = \begin{pmatrix} 1 & a & 2 \\ 0 & 1 & 1 \\ -1 & 1 & 1 \end{pmatrix}$,求 $a$ 及满足 $AP=B$ 的可逆矩阵 $P$.
	
	\item 已知 $x^3 - 5x + k = 0$ 有 3 个不同解,求 $k$ 取值范围.
	
	\item 设 $A$ 为 4 阶矩阵,$A^*$ 为其伴随矩阵.若 $Ax=0$ 的基础解系只有 2 个向量,求 $rank(A^*)$.
	
	\item 设 $A = \begin{pmatrix} 1 & 0 & -1 \\ 1 & 1 & -1 \\ 0 & 1 & a^2-1 \end{pmatrix}, b = \begin{pmatrix} 0 \\ 1 \\ a \end{pmatrix}$.若 $Ax=b$ 有无穷多解,求 $a$.
	
	\item 设 $I: \alpha_1 = \begin{pmatrix} 1 \\ 1 \\ 4 \end{pmatrix}, \alpha_2 = \begin{pmatrix} 1 \\ 0 \\ 4 \end{pmatrix}, \alpha_3 = \begin{pmatrix} 1 \\ 2 \\ a^2+3 \end{pmatrix}$；$II: \beta_1 = \begin{pmatrix} 1 \\ 1 \\ a+3 \end{pmatrix}, \beta_2 = \begin{pmatrix} 0 \\ 2 \\ 1-a \end{pmatrix}, \beta_3 = \begin{pmatrix} 1 \\ 3 \\ a^2+3 \end{pmatrix}$.若向量组 $I, II$ 等价,求 $a$,并将 $\beta_3$ 用 $\alpha_1, \alpha_2, \alpha_3$ 线性表示.
	
	\item 计算行列式:$D = \begin{vmatrix} a & 0 & -1 & 1 \\ 0 & a & 1 & -1 \\ -1 & 1 & a & 0 \\ 1 & -1 & 0 & a \end{vmatrix}$.
	
	\item 求 $f(x) = \begin{vmatrix} x & x & 1 & 2x \\ 1 & x & 2 & -1 \\ 2 & 1 & x & 1 \\ 2 & -1 & 1 & x \end{vmatrix}$ 中 $x^3$ 的系数.
	
	\item 设 $A = \begin{pmatrix} 1 & 1 & 1 \\ 1 & a & a^2 \\ 1 & b & b^2 \end{pmatrix}, b = \begin{pmatrix} 1 \\ 2 \\ 4 \end{pmatrix}$,判断 $Ax=b$ 解的情况.
	
	\item 设 $\alpha_1 = \begin{pmatrix} a \\ 1 \\ 1 \end{pmatrix}, \alpha_2 = \begin{pmatrix} 1 \\ a \\ 1 \end{pmatrix}, \alpha_3 = \begin{pmatrix} 1 \\ 1 \\ a \end{pmatrix}, \alpha_4 = \begin{pmatrix} 1 \\ a \\ a^2 \end{pmatrix}$.若 $\{\alpha_1, \alpha_2, \alpha_3\}$ 与 $\{\alpha_1, \alpha_2, \alpha_4\}$ 等价,求 $a$ 的取值范围.
	
	\item 设 $A$ 为 3 阶矩阵,交换 $A$ 的第 2 行和第 3 行,再将第 2 列的 $-1$ 倍加到第 1 列,得 $\begin{pmatrix} -2 & 1 & -1 \\ 1 & -1 & 0 \\ -1 & 0 & 0 \end{pmatrix}$,求 $A^{-1}$ 的迹.
	
	\item 设 $A, B$ 为 $n$ 阶可逆矩阵,求 $\begin{pmatrix} A & E \\ 0 & B \end{pmatrix}$ 的伴随矩阵.
	
	\item 设 $\alpha_1 = \begin{pmatrix} 1 \\ 2 \\ 3 \end{pmatrix}, \alpha_2 = \begin{pmatrix} 2 \\ 1 \\ 1 \end{pmatrix}, \beta_1 = \begin{pmatrix} 2 \\ 5 \\ 9 \end{pmatrix}, \beta_2 = \begin{pmatrix} 1 \\ 0 \\ 1 \end{pmatrix}$.若 $\gamma$ 可由 $\alpha_1, \alpha_2$ 线性表示,也可由 $\beta_1, \beta_2$ 线性表示,求 $\gamma$.
	
	\item 已知方程组 $\begin{cases} ax_1 + x_3 = 1 \\ x_1 + ax_2 + x_3 = 0 \\ x_1 + 2x_2 + ax_3 = 0 \\ ax_1 + bx_2 = 2 \end{cases}$ 有解.若 $\begin{vmatrix} a & 0 & 1 \\ 1 & a & 1 \\ 1 & 2 & a \end{vmatrix} = 4$,求 $\begin{vmatrix} 1 & a & 1 \\ 1 & 2 & a \\ a & b & 0 \end{vmatrix}$.
	
	\item 设 $A \begin{pmatrix} x_1 \\ x_2 \\ x_3 \end{pmatrix} = \begin{pmatrix} x_1 + x_2 + x_3 \\ 2x_1 - x_2 + x_3 \\ x_2 - x_3 \end{pmatrix}$,求 $A$ 及可逆矩阵 $P$ 和对角阵 $\Lambda$ 使得 $P^{-1}AP = \Lambda$.
	
	\item 设 $A$ 为 3 阶矩阵,$P = \begin{pmatrix} 1 & 0 & 0 \\ 0 & 1 & 0 \\ 1 & 0 & 1 \end{pmatrix}$.若 $P^T A P^2 = \begin{pmatrix} a+b & 0 & c \\ 0 & b & 0 \\ 2c & 0 & c \end{pmatrix}$,求 $A$.
	
	\item 设 $A = \begin{pmatrix} a+1 & b & 3 \\ a & \frac{b}{2} & 1 \\ 1 & 1 & 2 \end{pmatrix}$,$M_{ij}$ 是 $a_{ij}$ 的余子式.若 $|A| = -0.5$,且 $-M_{21} + M_{22} - M_{23} = 0$,求 $a, b$.
	
	\item 设 $A$ 为 3 阶矩阵,若 $r(2E-A) = 1, r(E+A) = 2$,求 $|A^{*}|$.
	
	\item 设 $A = \begin{pmatrix} 1 & -1 & 0 & -1 \\ 1 & 1 & 0 & 3 \\ 2 & 1 & 2 & 6 \end{pmatrix}, B = \begin{pmatrix} 1 & 0 & 1 & 2 \\ 1 & -1 & a & a - 1 \\ 2 & -3 & 2 & 2 \end{pmatrix}, \alpha = \begin{pmatrix} 0 \\ 2 \\ 3 \end{pmatrix}, \beta = \begin{pmatrix} 1 \\ 0 \\ -1 \end{pmatrix}$.
	\begin{enumerate}
		\item 证明 $Ax=\alpha$ 的解均为 $Bx=\beta$ 的解.
		\item 若 $Ax=\alpha$ 与 $Bx=\beta$ 不同解,求 $a$.
	\end{enumerate}
	

	
\end{enumerate}